\documentclass[12pt,a4paper]{article}
\usepackage[utf8]{inputenc}
\usepackage[T1]{fontenc}
\usepackage{amsmath,amssymb,amsthm}
\usepackage{physics}
\usepackage{geometry}
\usepackage{enumitem}
\usepackage{hyperref}
\usepackage{fancyhdr}
\usepackage{titlesec}

\geometry{margin=1in}
\pagestyle{fancy}
\fancyhf{}
\rhead{Lecture 7}
\lhead{Quantum Mechanics --- The Theoretical Minimum}
\cfoot{\thepage}

\titleformat{\section}{\large\bfseries}{}{0em}{}
\titleformat{\subsection}{\normalsize\bfseries}{}{0em}{}

\newtheorem{definition}{Definition}
\newtheorem{theorem}{Theorem}

\title{\textbf{Lecture 7: Entanglement and the Density Matrix} \\ \large Wave Functions, Reduced Descriptions, and Quantum Correlations}
\author{Based on Leonard Susskind's \textit{The Theoretical Minimum}}
\date{}

\begin{document}
\maketitle
\thispagestyle{fancy}

% ============================================================
\section{1. The Wave Function}
% ============================================================

\begin{definition}[Wave function]
Given a state $\ket{\Psi}$ and a complete orthonormal basis $\{\ket{a,b,c,\dots}\}$ labeled by eigenvalues of commuting observables, the \textbf{wave function} is
\[
    \psi(a,b,c,\dots) \equiv \braket{a,b,c,\dots}{\Psi}.
\]
It encodes the same information as $\ket{\Psi}$; every observable's expectation value can be computed from it.
\end{definition}

The state can be reconstructed from its wave function via
\[
    \ket{\Psi} = \sum_{a,b,c,\dots} \psi(a,b,c,\dots)\,\ket{a,b,c,\dots}.
\]

\subsection{1.1 Examples}
\begin{itemize}[nosep]
    \item \textbf{Single spin}: basis $\{\ket{\uparrow},\ket{\downarrow}\}$; wave function $\psi(\sigma_z)$ takes two values.
    \item \textbf{Two spins}: basis $\{\ket{\uparrow\uparrow},\ket{\uparrow\downarrow},\ket{\downarrow\uparrow},\ket{\downarrow\downarrow}\}$; wave function $\psi(\sigma_z,\tau_z)$ is a function of two discrete variables.
    \item \textbf{Particle on a line}: basis $\{\ket{x}\}$; wave function $\psi(x)$ is a continuous complex-valued function.
\end{itemize}

% ============================================================
\section{2. Expectation Values via the Wave Function}
% ============================================================

For an observable $L$ with matrix elements $L_{a'a}=\bra{a'}L\ket{a}$ (restricting to a single label for clarity):
\[
    \expval{L} = \sum_{a,a'} \psi^*(a')\,L_{a'a}\,\psi(a).
\]
This is the master formula from which all measurable predictions follow.

% ============================================================
\section{3. Composite Systems and the Wave Function}
% ============================================================

For a composite system of Alice ($a$) and Bob ($b$), the wave function is $\psi(a,b)$.

\begin{definition}[Product state]
The wave function factorizes:
\[
    \psi(a,b) = \psi_A(a)\,\phi_B(b).
\]
\end{definition}

An \textbf{entangled state} is one where no such factorization exists.

\subsection{3.1 Observables on composite systems}
Alice's observable $L$ has matrix elements
\[
    L_{a'b',\,ab} = L_{a'a}\,\delta_{b'b} \qquad\text{(identity on Bob)}.
\]
Bob's observable $M$:
\[
    M_{a'b',\,ab} = \delta_{a'a}\,M_{b'b} \qquad\text{(identity on Alice)}.
\]

% ============================================================
\section{4. The Density Matrix}
% ============================================================

When Alice has access only to her subsystem and cannot perform experiments on Bob, she needs a reduced description.

\begin{definition}[Alice's density matrix]
\[
    \boxed{\rho^{(A)}_{a,a'} = \sum_b \psi^*(a',b)\,\psi(a,b).}
\]
Bob's variables have been \textbf{summed over} (``traced out'').
\end{definition}

\subsection{4.1 Properties}
\begin{itemize}[nosep]
    \item $\rho^{(A)}$ is \textbf{Hermitian}: $(\rho^{(A)}_{a,a'})^* = \rho^{(A)}_{a',a}$.
    \item $\operatorname{Tr}\rho^{(A)}=1$ (normalization).
    \item For any Alice observable $L$: $\expval{L} = \operatorname{Tr}(L\,\rho^{(A)}) = \sum_{a,a'} L_{a'a}\,\rho^{(A)}_{a,a'}$.
    \item $\rho^{(A)}$ contains \emph{all} statistical information Alice can ever extract from her subsystem alone.
\end{itemize}

\subsection{4.2 Product state}
If $\psi(a,b)=\psi_A(a)\phi_B(b)$, then
\[
    \rho^{(A)}_{a,a'} = \psi_A(a)\,\psi_A^*(a') \cdot \underbrace{\sum_b |\phi_B(b)|^2}_{=1} = \psi_A(a)\,\psi_A^*(a').
\]
This is a \textbf{pure-state} density matrix: $(\rho^{(A)})^2 = \rho^{(A)}$.

\subsection{4.3 Entangled state}
For the singlet $\ket{\Psi}=\frac{1}{\sqrt{2}}(\ket{\uparrow\downarrow}-\ket{\downarrow\uparrow})$:
\[
    \rho^{(A)} = \frac{1}{2}\begin{pmatrix}1&0\\0&1\end{pmatrix} = \frac{1}{2}\mathbf{1}.
\]
This is a \textbf{maximally mixed} state: Alice's subsystem appears completely random.  She has \emph{zero} information about her spin alone, even though the composite state is a perfectly known pure state.

% ============================================================
\section{5. Correlations Revisited}
% ============================================================

\begin{definition}[Correlation]
\[
    C(L,M) = \expval{LM} - \expval{L}\expval{M}.
\]
\end{definition}

For the singlet:
\begin{align}
    \expval{\sigma_z} &= 0, & \expval{\tau_z} &= 0, & \expval{\sigma_z\tau_z} &= -1, \\
    \expval{\sigma_x} &= 0, & \expval{\tau_x} &= 0, & \expval{\sigma_x\tau_x} &= -1.
\end{align}
Every pair of matching components has correlation $-1$.  This is \textbf{maximal entanglement}: knowing one spin's measurement outcome determines the other's with certainty (opposite result).

\subsection{5.1 Criterion for entanglement}
A composite state is entangled if and only if there exists \emph{at least one} pair of subsystem observables $(L,M)$ with $C(L,M)\neq 0$.  Equivalently, the reduced density matrix $\rho^{(A)}$ is \emph{not} a pure state (i.e.\ $\operatorname{Tr}(\rho^{(A)})^2 < 1$).

% ============================================================
\section{6. Key Conceptual Points}
% ============================================================

\begin{enumerate}[nosep]
    \item A composite system can be in a definite pure state $\ket{\Psi}$ while each subsystem has \emph{no} definite state of its own---only a density matrix.
    \item This was Einstein, Podolsky, and Rosen's (EPR) central puzzle: ``We can know everything about the whole and nothing about the parts.''
    \item The density matrix formalism makes it precise: the subsystem's statistical description is obtained by tracing over the other subsystem's degrees of freedom.
    \item Entanglement is \emph{not} the same as ``sharing a state''---it is the failure of the wave function to factorize.
\end{enumerate}

% ============================================================
\section{7. Summary}
% ============================================================

\begin{enumerate}[nosep]
    \item The \textbf{wave function} $\psi(a,b,\dots)$ is the set of expansion coefficients of $\ket{\Psi}$ in a chosen basis.
    \item The \textbf{density matrix} $\rho^{(A)}$ describes a subsystem by tracing over the inaccessible degrees of freedom.
    \item In a product state, $\rho^{(A)}$ is pure; in an entangled state, it is mixed.
    \item Nonzero correlations between subsystem observables are the signature of entanglement.
    \item The singlet state is maximally entangled: every component measurement is perfectly anti-correlated, yet each individual spin appears maximally random.
\end{enumerate}

\end{document}
